


In this thesis, we investigated deep reservoir computing architectures over two different topological structures:
 cyclic and antisymmetric coupled modules.
This work extends the previous study on depth in recurrent neural networks and reservoir computing, 
addressing the problem of how depth can be effectively utilized in untrained recurrent systems.

Our aim was to introduce structured modular connectivity patterns that can help with information propagation and manipulation across multiple layers,
the inherent ability of deep reservoirs to not only capture sequential dependancies but also to develop a richer representation of different temporal scales,
made us ask if introducing a feedback mechanism across layers, such as we do in \DCR\@ could help with long-term memory and subsequentially to search for an architecture that
can achieve a longer memory retention in another way by introducing antisymmetric coupling between adjacent layers, as in \DAR.

From a theoretical standpoint, we derived the Jacobian of the global state transition map and established the necessary 
conditions for the Echo State Property (ESP).
Highlighting that in structured modular reservoirs, stability is a global property of the entire architecture rather than a layerwise condition.

Empirically, we evaluated the proposed architectures on memory capacity analysis, long temporal dependency benchmarks, and time-series tasks.
 Results show that structured modularity can enhance memory capacity and improve performance on tasks requiring long-term dependencies.
In particular, cyclic coupling reinforces global memory traces, while antisymmetric coupling helps with maintaining

Overall, we demonstrate that depth in reservoir computing can be exploited in a variety of ways, 
by carefully structuring a design in which topology directly affects how the model can process data and how information can flow, such that the memory of the model can be enhanced.


\section{Future Work}

Several research directions emerge from this study. 
\begin{itemize}
    \item Tighter theoretical characterizations of stability for block-cyclic and antisymmetric reservoirs could be developed beyond spectral radius bounds.
    \item Heterogeneous layer configurations may restore explicit time-scale differentiation while preserving structured feedback. 
    \item Mechanisms for controlled forgetting could mitigate excessive memory accumulation in strongly coupled modular systems. 
    \item Structured modular reservoirs offer promising directions for energy-efficient and hybrid reservoir computing systems.
\end{itemize}

The geometry of connectivity, therefore, emerges as a central design principle for deep reservoir architectures,
bridging dynamical systems theory and efficient sequence modeling.
